% Suggested insertion only. The original manuscript was not modified.
\subsection{Failure-Case Analysis and Practical Limitations}
\label{subsec:failure_case_analysis}

In addition to the aggregate Monte Carlo statistics, we examine a practically
relevant boundary failure in which mission-level target coverage is retained
but strict completion by every assigned interceptor is not. The trained
policies, target assignments, attacker trajectories, and the hit radius
$d_{\rm hit}=3~\mathrm{m}$ are kept fixed, and no policy update or retraining is
performed. To emulate imperfect battlefield information, the frozen-policy
evaluation introduces a one-sample ($0.05~\mathrm{s}$) observation delay and
zero-mean measurement noise with standard deviations of $3~\mathrm{m}$ and
$0.3~\mathrm{m/s}$ for position and velocity, respectively. We screen 100
reproducible episodes for each case. The reported examples are the earliest
three seeds per case satisfying the following fixed criterion: all eight
attackers are intercepted, exactly 19 of the 20 assigned interceptors hit,
seven of the eight assignment groups are complete, and the observed hit-time
spread does not exceed the $0.5~\mathrm{s}$ synchronization tolerance.

Fig.~\ref{fig:failure_cases} shows the resulting six boundary cases. In every
episode, assignment redundancy preserves complete target coverage
($8/8$), although one assigned interceptor fails to enter the lethal region
within the $75~\mathrm{s}$ evaluation horizon. As summarized in
Table~\ref{tab:failure_cases}, the missed closest-approach distances are only
$3.04$--$3.24~\mathrm{m}$ in Case~1 and $3.10$--$3.57~\mathrm{m}$ in Case~2.
Their margins above the hit threshold are therefore $0.04$--$0.24~\mathrm{m}$
and $0.10$--$0.57~\mathrm{m}$, respectively. These are near-threshold misses,
rather than grossly divergent engagements. The seven completed groups remain
well coordinated: their maximum observed hit-time spreads are
$0.05~\mathrm{s}$ in Case~1 and $0.15$--$0.35~\mathrm{s}$ in Case~2, all below
the prescribed tolerance.

The result distinguishes mission-level success from strict assignment-level
success. Redundant allocation localizes a single interceptor failure and still
neutralizes all attackers. However, the incomplete group's completion time is
right-censored beyond the evaluation horizon; hence the same episode represents
an unsuccessful individual interception and a delayed cooperative completion.
This exposes a limitation of the current framework: after an interceptor passes
its closest-approach point under delayed/noisy observations, the controller has
no explicit online miss-recovery or reassignment mechanism. Practical
extensions therefore include closing-rate-based miss detection,
event-triggered reassignment, reserve-interceptor activation, and
uncertainty-aware safety constraints. These examples are diagnostic boundary
cases and should not be interpreted as an estimate of the overall failure
probability.

\begin{figure*}[!t]
  \centering
  \includegraphics[width=\textwidth]{figures/failure_cases_case1_case2_v10.pdf}
  \caption{Representative boundary failures under frozen-policy evaluation.
  Solid curves denote attacker trajectories, dotted curves successful
  interceptors, and the thick dashed curve the missed assigned interceptor.
  Green circles mark successful hit events and the cross marks the missed
  interceptor's closest-approach point. In all six episodes, 19/20 assigned
  interceptors hit and all 8/8 attackers are nevertheless intercepted.}
  \label{fig:failure_cases}
\end{figure*}

\begin{table*}[!t]
\centering
\caption{Quantitative failure-case diagnostics.}
\label{tab:failure_cases}
\begin{tabular}{cccccccc}
\hline
Case & Seed & Hits & Targets & Missed ID & Assigned target &
$d_{\min}^{\rm miss}$ (m) & Max. observed spread (s) \\
\hline
1 & 76014 & 19/20 & 8/8 & 1  & 21 & 3.161 & 0.050 \\
1 & 76048 & 19/20 & 8/8 & 5  & 25 & 3.243 & 0.050 \\
1 & 76052 & 19/20 & 8/8 & 8  & 20 & 3.041 & 0.050 \\
2 & 77008 & 19/20 & 8/8 & 1  & 21 & 3.574 & 0.350 \\
2 & 77020 & 19/20 & 8/8 & 2  & 22 & 3.096 & 0.150 \\
2 & 77023 & 19/20 & 8/8 & 17 & 21 & 3.488 & 0.200 \\
\hline
\end{tabular}
\end{table*}

